\chapter{Modelo de Datos}

Esta sección describe el modelo tal como existe en \texttt{db/lab1\_desigualdad.sqlite}. No se corrige ni reinterpreta el diseño observado; se documenta su implementación real. El modelo responde a una lógica dimensión-hecho complementada por una tabla de metadata.

\begin{table}[htbp]
\centering
\footnotesize
\caption{Resumen del modelo físico observado en SQLite.}
\begin{tabular}{|p{3.5cm}|p{4.2cm}|p{2.2cm}|r|r|p{3.2cm}|}
\hline
Tabla & Rol & Llave primaria & Cols. & Filas & Observación \\
\hline
dim\_comuna & Dimensión de comunas del universo final. & codigo\_comuna & 5 & 32 & Conserva nombre, provincia, región y fuente\_referencia. \\
fact\_desigualdad\_comunal & Tabla de hechos con variables finales y métricas derivadas. & codigo\_comuna & 11 & 32 & No replica nombre\_comuna; se relaciona con dim\_comuna. \\
metadata\_fuentes & Contrato operativo y trazabilidad de lectura de fuentes. & id\_fuente & 13 & 4 & Documenta institución, URL, hoja, skiprows y archivo lógico. \\
\hline
\end{tabular}
\end{table}


\section{Tabla \texttt{dim\_comuna}}

La dimensión comunal contiene 32 filas y cinco columnas: \texttt{codigo\_comuna}, \texttt{nombre\_comuna}, \texttt{provincia}, \texttt{region} y \texttt{fuente\_referencia}. Esta tabla preserva la identificación y el contexto territorial del universo final.

\section{Tabla \texttt{fact\_desigualdad\_comunal}}

La tabla de hechos contiene 32 filas y once columnas. Su llave es \texttt{codigo\_comuna} y almacena \texttt{poblacion}, \texttt{anio\_poblacion}, \texttt{pobreza\_ingresos\_pct}, \texttt{anio\_pobreza}, \texttt{areas\_verdes\_m2}, \texttt{anio\_areas\_verdes}, \texttt{ipp\_miles\_pesos}, \texttt{anio\_ingresos}, \texttt{areas\_verdes\_m2\_hab} e \texttt{ipp\_pesos\_hab}. El campo \texttt{nombre\_comuna} no se replica en esta tabla; se obtiene mediante unión con \texttt{dim\_comuna}.

\section{Tabla \texttt{metadata\_fuentes}}

La tabla \texttt{metadata\_fuentes} contiene 4 filas y trece columnas. Su función es conservar la trazabilidad operativa de las fuentes: institución, URL, fecha de descarga, formato, año de referencia, hoja, \texttt{skiprows}, archivo lógico y observaciones. Esta tabla conecta el resultado analítico con el contrato de lectura real del ETL.

\section{Relación entre tablas}

El vínculo entre la dimensión y la tabla de hechos se realiza por \texttt{codigo\_comuna}. En términos operativos, la combinación \texttt{dim\_comuna} + \texttt{fact\_desigualdad\_comunal} reproduce el dataset final validado, mientras que \texttt{metadata\_fuentes} documenta el origen de los insumos y las reglas mínimas de lectura.
